\documentclass[10pt]{article}
\usepackage[T1]{fontenc}
\usepackage[utf8]{inputenc}
\usepackage[french]{babel}
\usepackage{amsmath,amssymb,mathrsfs,mathpazo}
\usepackage{geometry}
\geometry{a4paper,margin=2cm}
\usepackage{graphicx,array,tabularx,tabularray,colortbl,multirow,multicol,booktabs,enumitem,pifont,ulem,fancybox,diagbox,cancel}
\usepackage{pstricks,pstricks-add,pst-plot,pst-tree,pst-node,pst-text,pst-func,pst-math,pst-bspline,pst-3dplot,pst-eucl,pst-solides3d}
\let\pscalculate\undefined
\usepackage{pst-bezier}
\usepackage{tikz,pgfplots,tkz-tab}
\usetikzlibrary{arrows,arrows.meta,patterns,calc,positioning,trees,decorations.markings,shapes.geometric,shapes.misc,shapes.arrows,intersections,decorations.pathmorphing,decorations.pathreplacing,backgrounds,shadows,babel}
\pgfplotsset{compat=1.18}
\usepgfplotslibrary{fillbetween}
\usepackage[np]{numprint}
\usepackage{listings,makecell,pas-tableur,twoopt}
\let\barre\undefined
\usepackage[pstricks]{tablvar}
\definecolor{darkgreen}{RGB}{0,100,0}
\providecommand{\red}{\color{red}}
\providecommand{\blue}{\color{blue}}
\providecommand{\green}{\color{green}}
\providecommand{\black}{\color{black}}
\providecommand{\white}{\color{white}}
\providecommand{\gray}{\color{gray}}
\providecommand{\cyan}{\color{cyan}}
\providecommand{\magenta}{\color{magenta}}
\providecommand{\yellow}{\color{yellow}}
\providecommand{\violet}{\color{violet}}
\providecommand{\orange}{\color{orange}}
\providecommand{\euro}{\texteuro}
\providecommand{\EUR}{\texteuro}
\providecommand{\ieme}{\textsuperscript{e}}
\providecommand{\degres}{\ensuremath{^\circ}}
\providecommand{\degree}{\ensuremath{^\circ}}
\providecommand{\vv}[1]{\overrightarrow{#1}}
\providecommand{\Coord}[1]{#1}
\providecommand{\decosix}{}
\providecommand{\starredbullet}{\ensuremath{\bullet}}
\pagestyle{empty}
\setlength{\parindent}{0pt}
\usepackage{adjustbox}\DeclareUnicodeCharacter{2236}{:}\DeclareUnicodeCharacter{2190}{\ensuremath{\leftarrow}}
\providecommand{\R}{\mathbb{R}}
\providecommand{\N}{\mathbb{N}}
\providecommand{\Z}{\mathbb{Z}}
\providecommand{\Q}{\mathbb{Q}}
\providecommand{\D}{\mathbb{D}}
\providecommand{\C}{\mathbb{C}}
\providecommand{\e}{{\mathrm{e}}}
\providecommand{\vect}[1]{\overrightarrow{#1}}
\providecommand{\vectt}[1]{\overrightarrow{\mathrm{#1}}}
\providecommand{\vvt}[1]{\overrightarrow{\mathrm{#1}}}
\providecommand{\np}[2][]{{#2~#1}}
\providecommand{\vv}[1]{{\overrightarrow{#1}}}
\renewcommand{\d}{{\mathrm{d}}}
\providecommand{\degres}{{^\circ}}
\providecommand{\degre}{{^\circ}}
\providecommand{\blue}{\color{blue}}
\providecommand{\red}{\color{red}}
\providecommand{\enskip}{\quad}
\providecommand{\llbracket}{\lbrack\!\lbrack}
\providecommand{\rrbracket}{\rbrack\!\rbrack}
\providecommand{\euro}{{\text{€}}}
\providecommand{\no}{{\text{n°}}}
\providecommand{\interval}[2]{{[#1;#2]}}
\providecommand{\intervFO}[2]{{[#1;#2[}}
\providecommand{\dotfill}{\ldots}

\providecommand{\ds}{\displaystyle}
\providecommand{\textding}[1]{#1}
\providecommand{\Oij}{\(\mathrm{O};\vec{\imath},\vec{\jmath}\)}
\providecommand{\Oijk}{\(\mathrm{O};\vec{\imath},\vec{\jmath},\vec{k}\)}
\providecommand{\Ouv}{\(\mathrm{O};\vec{u},\vec{v}\)}
\providecommand{\intervFF}[2]{\ensuremath{[#1;#2]}}
\renewcommand{\euro}{\ifmmode\text{\texteuro}\else\texteuro\fi}
\DeclareUnicodeCharacter{00B0}{\ifmmode^{\circ}\else\textdegree\fi}
\DeclareUnicodeCharacter{20AC}{\ifmmode\text{\texteuro}\else\texteuro\fi}
% Énoncés : Éducation nationale. Transcription et corrigés : APMEP, Denis Vergès et contributeurs. Découpage et classement : CoopMaths, Sébastien Lozano, Éric Elter et contributeurs. Adaptation de la présentation : Maths Explorer.
% CC BY-SA : https://coopmaths.fr/www/about/
% https://coopmaths.fr/alea/static/eam/2025/tex/eam_technologique_2025_06_sujet0v1_3.tex
\begin{document}
\noindent\textit{APMEP / CoopMaths. CC BY-SA. Présentation adaptée par Maths Explorer.}\par\medskip
\begin{minipage}[t]{0.6\textwidth}
On considère une fonction $f$ définie et dérivable sur l'intervalle $[-2~;~6]$. \\
Sa courbe représentative, notée $C$, est donnée ci-contre.

\begin{itemize}
\item On sait que la courbe $C$ passe par les points de coordonnées $(0~;~8)$, $(2~;~0)$ et $(4~;~-8)$.
\item On note $T$ la tangente à la courbe $C$ au point d'abscisse $x = 2$.
\item On sait que la tangente $T$ coupe l'axe des ordonnées en $y = 12$.
\end{itemize}

On note $f'$ la fonction dérivée de $f$.
\end{minipage}
\hfill
\begin{minipage}[t]{0.35\textwidth}
\vspace{0pt}
\begin{center}
\begin{tikzpicture}[scale=0.75, transform shape]
  \begin{axis}[
    xmin=-2.5, xmax=6.5,
    ymin=-15, ymax=18,
    samples=100,
    axis lines=middle,
    xlabel=\(x\),
    ylabel=\(y\),
    ytick={5},
    xtick={1},
    extra y ticks={-15,-10,-5,0,10,15},
    extra x ticks={-2,-1,0,2,3,4,5,6},
    extra tick style={tick label style={opacity=0}},
    grid=both,
    grid style={line width=.1pt, draw=gray!10},
    major grid style={line width=.2pt,draw=gray!50},
    axis line style={line width=1pt},
    xtick distance=1,
    ytick distance=5,
    xmajorgrids,
    ymajorgrids,
    restrict y to domain=-15:18,
    restrict x to domain=-2.5:6.5,
  ]
  \node[below left] at (axis cs:0,0) {0};
  \node[blue, below right] at (axis cs:5,-5) {C};
  \addplot[blue, very thick, domain=-2:6] {0.5*x^3 - 3*x^2 + 8};
  \node[red, below left] at (axis cs:-0.5,15) {T};
  \addplot[red, very thick, domain=-1:4.5] {-6*x + 12};
  \end{axis}
\end{tikzpicture}
\end{center}
\end{minipage}

\begin{enumerate}
\item 
\begin{enumerate}
\item Déterminer les valeurs de $f(2)$ et $f'(2)$.
\item Donner une équation de la tangente $T$.
\item Recopier et compléter le tableau de variation ci-dessous en utilisant le graphique.

\begin{center}
\setlength{\tabcolsep}{18pt}
\begin{tabular}{|c|cccc|}
\hline
$x$ & $-2$ & $0$ & $4$ & $6$ \\
\hline
 & & & & \\
Variations & & & & \\
de $f$ & & & & \\
 & & & & \\
\hline
\end{tabular}
\end{center}
\end{enumerate}

\item On admet que la fonction $f$ est définie sur l'intervalle $[-2~;~6]$ par $f(x) = 0,5x^3 - 3x^2 + 8$.
\begin{enumerate}
\item Montrer que, pour tout réel $x$ de l'intervalle $[-2~;~6]$, on a $f'(x) = 1,5x(x - 4)$.
\item Étudier le signe de $f'(x)$ et retrouver le tableau de variation de la fonction $f$ sur l'intervalle $[-2~;~6]$.
\end{enumerate}

\item On admet que, pour tout réel $x$ de l'intervalle $[0~;~2]$, on a $f(x) \leq -6x + 12$. \\
Que peut-on en déduire pour la courbe $C$ et la tangente $T$ sur l'intervalle $[0~;~2]$ ?
\end{enumerate}

\end{document}
